\chapter{Uvod}

Natjecateljsko programiranje spaja dvije teme:
(1) osmišljavanje algoritama i (2) implementaciju algoritama.

\key{Osmišljavanje algoritama} sastoji se od rješavanja zadataka
i matematičkog razmišljanja.
Potrebne su vještine analiziranja zadataka i njihova
kreativnog rješavanja.
Algoritam za rješavanje zadatka
mora biti i točan i efikasan,
a srž zadatka često je
u osmišljavanju efikasnog algoritma.

Teorijsko znanje o algoritmima
važno je natjecateljima u programiranju.
Rješenje zadatka tipično je
kombinacija dobro poznatih tehnika i
novih uvida.
Tehnike koje se pojavljuju u natjecateljskom programiranju
ujedno čine temelj znanstvenog istraživanja
algoritama.

\key{Implementacija algoritama} zahtijeva dobre
programerske vještine.
U natjecateljskom programiranju rješenja se
ocjenjuju testiranjem implementiranog algoritma
na skupu testnih primjera.
Stoga nije dovoljno da je ideja
algoritma točna, nego i implementacija
mora biti točna.

Dobar stil pisanja koda na natjecanjima
jednostavan je i sažet.
Programe treba pisati brzo,
jer nema puno raspoloživog vremena.
Za razliku od tradicionalnog programskog inženjerstva,
programi su kratki (obično najviše nekoliko
stotina redaka koda) i ne moraju se 
održavati nakon natjecanja.

\section{Programski jezici}

\index{programski jezik}

Trenutno su najpopularniji programski
jezici koji se koriste na natjecanjima C++, Python i Java.
Na primjer, na Google Code Jamu 2017.,
među 3.000 najboljih sudionika,
79 \% koristilo je C++,
16 \% Python i
8 \% Javu \cite{goo17}.
Neki sudionici koristili su i više jezika.

Mnogi smatraju da je C++ najbolji izbor
za natjecatelja u programiranju,
a C++ je gotovo uvijek dostupan u
natjecateljskim sustavima.
Prednosti korištenja C++-a su te da
je to vrlo efikasan jezik i
da njegova standardna biblioteka sadrži 
veliku zbirku
struktura podataka i algoritama.

S druge strane, dobro je
vladati s više jezika i razumjeti
njihove prednosti.
Na primjer, ako su u zadatku potrebni
veliki cijeli brojevi,
Python može biti dobar izbor jer
sadrži ugrađene operacije za
računanje s velikim cijelim brojevima.
Ipak, većina zadataka na programerskim natjecanjima
postavljena je tako da
korištenje određenog programskog jezika
nije nepravedna prednost.

Svi primjeri programa u ovoj knjizi napisani su u C++-u,
a često se koriste strukture podataka i algoritmi
iz standardne biblioteke.
Programi slijede standard C++11,
koji se danas može koristiti na većini natjecanja.
Ako još ne znate programirati u C++-u,
sada je dobar trenutak da počnete učiti.

\subsubsection{Predložak C++ koda}

Tipičan predložak C++ koda za natjecateljsko programiranje
izgleda ovako:

\begin{lstlisting}
#include <bits/stdc++.h>

using namespace std;

int main() {
    // solution comes here
}
\end{lstlisting}

Redak \texttt{\#include} na početku
koda značajka je prevoditelja (compiler) \texttt{g++}
koja nam omogućuje uključivanje cijele standardne biblioteke.
Stoga nije potrebno zasebno uključivati
biblioteke poput \texttt{iostream},
\texttt{vector} i \texttt{algorithm},
nego su one automatski dostupne.

Redak \texttt{using} deklarira
da se klase i funkcije
standardne biblioteke mogu izravno koristiti
u kodu.
Bez retka \texttt{using} morali bismo
pisati, primjerice, \texttt{std::cout},
a sada je dovoljno napisati \texttt{cout}.

Kod se može prevesti sljedećom naredbom:

\begin{lstlisting}
g++ -std=c++11 -O2 -Wall test.cpp -o test
\end{lstlisting}

Ova naredba iz izvornog koda \texttt{test.cpp}
stvara binarnu datoteku \texttt{test}.
Prevoditelj slijedi standard C++11
(\texttt{-std=c++11}),
optimizira kod (\texttt{-O2})
i prikazuje upozorenja o mogućim pogreškama (\texttt{-Wall}).

\section{Ulaz i izlaz}

\index{ulaz i izlaz}

Na većini natjecanja za čitanje ulaza i ispis izlaza
koriste se standardni tokovi.
U C++-u standardni su tokovi
\texttt{cin} za ulaz i \texttt{cout} za izlaz.
Uz to se mogu koristiti i C-ove funkcije
\texttt{scanf} i \texttt{printf}.

Ulaz programa obično se sastoji od
brojeva i znakovnih nizova odvojenih
razmacima i novim redovima.
Mogu se čitati iz toka \texttt{cin}
na sljedeći način:

\begin{lstlisting}
int a, b;
string x;
cin >> a >> b >> x;
\end{lstlisting}

Ovakav kod uvijek radi,
uz pretpostavku da je između svakog elementa u ulazu
barem jedan razmak ili novi red.
Na primjer, gornji kod može pročitati
oba sljedeća ulaza:
\begin{lstlisting}
123 456 monkey
\end{lstlisting}
\begin{lstlisting}
123    456
monkey
\end{lstlisting}
Tok \texttt{cout} koristi se za izlaz
na sljedeći način:
\begin{lstlisting}
int a = 123, b = 456;
string x = "monkey";
cout << a << " " << b << " " << x << "\n";
\end{lstlisting}

Ulaz i izlaz ponekad su
usko grlo programa.
Sljedeći redci na početku koda
čine ulaz i izlaz efikasnijima:

\begin{lstlisting}
ios::sync_with_stdio(0);
cin.tie(0);
\end{lstlisting}

Uočimo da novi red \texttt{"\textbackslash n"}
radi brže od \texttt{endl},
jer \texttt{endl} uvijek uzrokuje
operaciju pražnjenja spremnika.

C-ove funkcije \texttt{scanf}
i \texttt{printf} alternativa su
standardnim tokovima C++-a.
Obično su nešto brže,
ali su i teže za korištenje.
Sljedeći kod čita dva cijela broja s ulaza:
\begin{lstlisting}
int a, b;
scanf("%d %d", &a, &b);
\end{lstlisting}
Sljedeći kod ispisuje dva cijela broja:
\begin{lstlisting}
int a = 123, b = 456;
printf("%d %d\n", a, b);
\end{lstlisting}

Ponekad program treba pročitati cijeli redak
s ulaza, koji možda sadrži razmake.
To se može postići korištenjem
funkcije \texttt{getline}:

\begin{lstlisting}
string s;
getline(cin, s);
\end{lstlisting}

Ako količina podataka nije poznata, korisna je
sljedeća petlja:
\begin{lstlisting}
while (cin >> x) {
    // code
}
\end{lstlisting}
Ova petlja čita elemente s ulaza
jedan za drugim, sve dok na ulazu
više nema dostupnih podataka.

U nekim natjecateljskim sustavima za ulaz i izlaz
koriste se datoteke.
Jednostavno rješenje za to jest napisati
kod kao i obično koristeći standardne tokove,
ali dodati sljedeće retke na početak koda:
\begin{lstlisting}
freopen("input.txt", "r", stdin);
freopen("output.txt", "w", stdout);
\end{lstlisting}
Nakon toga program čita ulaz iz datoteke
''input.txt'' i piše izlaz u datoteku
''output.txt''.

\section{Rad s brojevima}

\index{cijeli broj}

\subsubsection{Cijeli brojevi}

Najčešće korišten cjelobrojni tip u natjecateljskom programiranju
jest \texttt{int}, koji je 32-bitni tip s
rasponom vrijednosti $-2^{31} \ldots 2^{31}-1$
ili otprilike $-2 \cdot 10^9 \ldots 2 \cdot 10^9$.
Ako tip \texttt{int} nije dovoljan,
može se koristiti 64-bitni tip \texttt{long long}.
On ima raspon vrijednosti $-2^{63} \ldots 2^{63}-1$
ili otprilike $-9 \cdot 10^{18} \ldots 9 \cdot 10^{18}$.

Sljedeći kod definira
varijablu tipa \texttt{long long}:
\begin{lstlisting}
long long x = 123456789123456789LL;
\end{lstlisting}
Sufiks \texttt{LL} znači da je
tip broja \texttt{long long}.

Česta pogreška pri korištenju tipa \texttt{long long}
jest da se negdje u kodu i dalje koristi
tip \texttt{int}.
Na primjer, sljedeći kod sadrži
suptilnu pogrešku:

\begin{lstlisting}
int a = 123456789;
long long b = a*a;
cout << b << "\n"; // -1757895751
\end{lstlisting}

Iako je varijabla \texttt{b} tipa \texttt{long long},
oba broja u izrazu \texttt{a*a}
tipa su \texttt{int}, pa je i rezultat
tipa \texttt{int}.
Zbog toga će varijabla \texttt{b}
sadržavati pogrešan rezultat.
Zadatak se može riješiti promjenom tipa
varijable \texttt{a} u \texttt{long long} ili
promjenom izraza u \texttt{(long long)a*a}.

Natjecateljski zadaci obično su postavljeni tako da je
tip \texttt{long long} dovoljan.
Ipak, dobro je znati da
prevoditelj \texttt{g++} nudi i
128-bitni tip \texttt{\_\_int128\_t}
s rasponom vrijednosti
$-2^{127} \ldots 2^{127}-1$ ili otprilike $-10^{38} \ldots 10^{38}$.
Međutim, taj tip nije dostupan u svim natjecateljskim sustavima.

\subsubsection{Modularna aritmetika}

\index{ostatak}
\index{modularna aritmetika}

S $x \bmod m$ označavamo ostatak
pri dijeljenju $x$ s $m$.
Na primjer, $17 \bmod 5 = 2$,
jer je $17 = 3 \cdot 5 + 2$.

Ponekad je odgovor na zadatak
vrlo velik broj, ali ga je dovoljno
ispisati ''modulo $m$'', tj.
kao ostatak pri dijeljenju odgovora s $m$
(na primjer, ''modulo $10^9+7$'').
Ideja je da je, čak i ako je stvarni odgovor
vrlo velik,
dovoljno koristiti tipove
\texttt{int} i \texttt{long long}.

Važno svojstvo ostatka jest da se
kod zbrajanja, oduzimanja i množenja
ostatak može uzeti prije operacije:

\[
\begin{array}{rcr}
(a+b) \bmod m & = & (a \bmod m + b \bmod m) \bmod m \\
(a-b) \bmod m & = & (a \bmod m - b \bmod m) \bmod m \\
(a \cdot b) \bmod m & = & (a \bmod m \cdot b \bmod m) \bmod m
\end{array}
\]

Dakle, možemo uzeti ostatak nakon svake operacije
i brojevi nikada neće postati preveliki.

Na primjer, sljedeći kod računa $n!$,
faktorijelu broja $n$, modulo $m$:
\begin{lstlisting}
long long x = 1;
for (int i = 2; i <= n; i++) {
    x = (x*i)%m;
}
cout << x%m << "\n";
\end{lstlisting}

Obično želimo da ostatak uvijek
bude između $0\ldots m-1$.
Međutim, u C++-u i drugim jezicima
ostatak negativnog broja
jednak je nuli ili je negativan.
Jednostavan način da osiguramo da
nema negativnih ostataka jest da najprije izračunamo
ostatak kao i obično, a zatim dodamo $m$
ako je rezultat negativan:
\begin{lstlisting}
x = x%m;
if (x < 0) x += m;
\end{lstlisting}
Međutim, to je potrebno samo kada u
kodu postoje oduzimanja pa
ostatak može postati negativan.

\subsubsection{Brojevi s pomičnim zarezom}

\index{broj s pomičnim zarezom}

Uobičajeni tipovi s pomičnim zarezom u
natjecateljskom programiranju su
64-bitni \texttt{double}
i, kao proširenje prevoditelja \texttt{g++},
80-bitni \texttt{long double}.
U većini slučajeva \texttt{double} je dovoljan,
ali \texttt{long double} je precizniji.

Potrebna preciznost odgovora
obično je zadana u tekstu zadatka.
Jednostavan način ispisa odgovora jest korištenje
funkcije \texttt{printf}
i navođenje broja decimalnih mjesta
u formatnom nizu.
Na primjer, sljedeći kod ispisuje
vrijednost $x$ s 9 decimalnih mjesta:

\begin{lstlisting}
printf("%.9f\n", x);
\end{lstlisting}

Poteškoća pri korištenju brojeva s pomičnim zarezom
jest da se neki brojevi ne mogu točno
prikazati kao brojevi s pomičnim zarezom,
pa će doći do pogrešaka zaokruživanja.
Na primjer, rezultat sljedećeg koda
je iznenađujući:

\begin{lstlisting}
double x = 0.3*3+0.1;
printf("%.20f\n", x); // 0.99999999999999988898
\end{lstlisting}

Zbog pogreške zaokruživanja
vrijednost varijable \texttt{x} malo je manja od 1,
dok bi točna vrijednost bila 1.

Riskantno je uspoređivati brojeve s pomičnim zarezom
operatorom \texttt{==},
jer je moguće da bi vrijednosti trebale biti
jednake, ali nisu zbog pogrešaka u preciznosti.
Bolji način usporedbe brojeva s pomičnim zarezom
jest pretpostaviti da su dva broja jednaka
ako je razlika među njima manja od $\varepsilon$,
gdje je $\varepsilon$ malen broj.

U praksi se brojevi mogu usporediti
ovako ($\varepsilon=10^{-9}$):

\begin{lstlisting}
if (abs(a-b) < 1e-9) {
    // a and b are equal
}
\end{lstlisting}

Uočimo da se, iako su brojevi s pomičnim zarezom netočni,
cijeli brojevi do određene granice ipak mogu
prikazati točno.
Na primjer, koristeći \texttt{double},
moguće je točno prikazati sve
cijele brojeve čija je apsolutna vrijednost najviše $2^{53}$.

\section{Skraćivanje koda}

Kratak kod idealan je u natjecateljskom programiranju,
jer programe treba pisati
što je brže moguće.
Zbog toga natjecatelji često definiraju
kraća imena za tipove podataka i druge dijelove koda.

\subsubsection{Imena tipova}
\index{tuppdef@\texttt{typedef}}
Naredbom \texttt{typedef}
moguće je tipu podataka dati
kraće ime.
Na primjer, ime \texttt{long long} je dugo,
pa možemo definirati kraće ime \texttt{ll}:
\begin{lstlisting}
typedef long long ll;
\end{lstlisting}
Nakon toga kod
\begin{lstlisting}
long long a = 123456789;
long long b = 987654321;
cout << a*b << "\n";
\end{lstlisting}
može se skratiti ovako:
\begin{lstlisting}
ll a = 123456789;
ll b = 987654321;
cout << a*b << "\n";
\end{lstlisting}

Naredba \texttt{typedef}
može se koristiti i za složenije tipove.
Na primjer, sljedeći kod daje
ime \texttt{vi} vektoru cijelih brojeva
i ime \texttt{pi} paru
koji sadrži dva cijela broja.
\begin{lstlisting}
typedef vector<int> vi;
typedef pair<int,int> pi;
\end{lstlisting}

\subsubsection{Makroi}
\index{makro}
Drugi način skraćivanja koda jest definiranje
\key{makroa}.
Makro znači da će se određeni znakovni nizovi u
kodu promijeniti prije prevođenja.
U C++-u makroi se definiraju
ključnom riječi \texttt{\#define}.

Na primjer, možemo definirati sljedeće makroe:
\begin{lstlisting}
#define F first
#define S second
#define PB push_back
#define MP make_pair
\end{lstlisting}
Nakon toga kod
\begin{lstlisting}
v.push_back(make_pair(y1,x1));
v.push_back(make_pair(y2,x2));
int d = v[i].first+v[i].second;
\end{lstlisting}
može se skratiti ovako:
\begin{lstlisting}
v.PB(MP(y1,x1));
v.PB(MP(y2,x2));
int d = v[i].F+v[i].S;
\end{lstlisting}

Makro može imati i parametre,
što omogućuje skraćivanje petlji i drugih
struktura.
Na primjer, možemo definirati sljedeći makro:
\begin{lstlisting}
#define REP(i,a,b) for (int i = a; i <= b; i++)
\end{lstlisting}
Nakon toga kod
\begin{lstlisting}
for (int i = 1; i <= n; i++) {
    search(i);
}
\end{lstlisting}
može se skratiti ovako:
\begin{lstlisting}
REP(i,1,n) {
    search(i);
}
\end{lstlisting}

Makroi ponekad uzrokuju greške koje je teško
otkriti. Na primjer, promotrimo sljedeći makro
koji računa kvadrat broja:
\begin{lstlisting}
#define SQ(a) a*a
\end{lstlisting}
Ovaj makro \emph{ne radi} uvijek onako kako očekujemo.
Na primjer, kod
\begin{lstlisting}
cout << SQ(3+3) << "\n";
\end{lstlisting}
odgovara kodu
\begin{lstlisting}
cout << 3+3*3+3 << "\n"; // 15
\end{lstlisting}

Bolja inačica makroa je sljedeća:
\begin{lstlisting}
#define SQ(a) (a)*(a)
\end{lstlisting}
Sada kod
\begin{lstlisting}
cout << SQ(3+3) << "\n";
\end{lstlisting}
odgovara kodu
\begin{lstlisting}
cout << (3+3)*(3+3) << "\n"; // 36
\end{lstlisting}


\section{Matematika}

Matematika igra važnu ulogu u natjecateljskom
programiranju i nije moguće postati
uspješan natjecatelj u programiranju bez
dobrih matematičkih vještina.
Ovo poglavlje razmatra neke važne
matematičke pojmove i formule koje
su nam potrebne kasnije u knjizi.

\subsubsection{Formule za sume}

Svaka suma oblika
\[\sum_{x=1}^n x^k = 1^k+2^k+3^k+\ldots+n^k,\]
gdje je $k$ pozitivan cijeli broj,
ima zatvorenu formulu koja je
polinom stupnja $k+1$.
Na primjer\footnote{\index{Faulhaberova formula}
Postoji čak i opća formula za takve sume, zvana \key{Faulhaberova formula},
ali je previše složena da bismo je ovdje prikazali.},
\[\sum_{x=1}^n x = 1+2+3+\ldots+n = \frac{n(n+1)}{2}\]
i
\[\sum_{x=1}^n x^2 = 1^2+2^2+3^2+\ldots+n^2 = \frac{n(n+1)(2n+1)}{6}.\]

\key{Aritmetički niz} je \index{aritmetički niz}
niz brojeva
u kojem je razlika između svaka dva uzastopna
broja konstantna.
Na primjer,
\[3, 7, 11, 15\]
je aritmetički niz s konstantom 4.
Suma aritmetičkog niza može se izračunati
formulom
\[\underbrace{a + \cdots + b}_{n \,\, \textrm{brojeva}} = \frac{n(a+b)}{2}\]
gdje je $a$ prvi broj,
$b$ posljednji broj, a
$n$ količina brojeva.
Na primjer,
\[3+7+11+15=\frac{4 \cdot (3+15)}{2} = 36.\]
Formula se temelji na tome
da se suma sastoji od $n$ brojeva i
da je vrijednost svakog broja u prosjeku $(a+b)/2$.

\index{geometrijski niz}
\key{Geometrijski niz} je niz
brojeva
u kojem je odnos između svaka dva uzastopna
broja konstantan.
Na primjer,
\[3,6,12,24\]
je geometrijski niz s konstantom 2.
Suma geometrijskog niza može se izračunati
formulom
\[a + ak + ak^2 + \cdots + b = \frac{bk-a}{k-1}\]
gdje je $a$ prvi broj,
$b$ posljednji broj, a
odnos između uzastopnih brojeva je $k$.
Na primjer,
\[3+6+12+24=\frac{24 \cdot 2 - 3}{2-1} = 45.\]

Ova formula može se izvesti na sljedeći način. Neka je
\[ S = a + ak + ak^2 + \cdots + b .\]
Množenjem obiju strana s $k$ dobivamo
\[ kS = ak + ak^2 + ak^3 + \cdots + bk,\]
a rješavanje jednadžbe
\[ kS-S = bk-a\]
daje formulu.

Poseban slučaj sume geometrijskog niza je formula
\[1+2+4+8+\ldots+2^{n-1}=2^n-1.\]

\index{harmonijska suma}

\key{Harmonijska suma} je suma oblika
\[ \sum_{x=1}^n \frac{1}{x} = 1+\frac{1}{2}+\frac{1}{3}+\ldots+\frac{1}{n}.\]

Gornja ograda harmonijske sume je $\log_2(n)+1$.
Naime, možemo
promijeniti svaki član $1/k$ tako da $k$ postane
najbliža potencija dvojke koja ne prelazi $k$.
Na primjer, kada je $n=6$, sumu možemo
procijeniti na sljedeći način:
\[ 1+\frac{1}{2}+\frac{1}{3}+\frac{1}{4}+\frac{1}{5}+\frac{1}{6} \le
1+\frac{1}{2}+\frac{1}{2}+\frac{1}{4}+\frac{1}{4}+\frac{1}{4}.\]
Ova gornja ograda sastoji se od $\log_2(n)+1$ dijelova
($1$, $2 \cdot 1/2$, $4 \cdot 1/4$, itd.),
a vrijednost svakog dijela je najviše 1.

\subsubsection{Teorija skupova}

\index{teorija skupova}
\index{skup}
\index{presjek}
\index{unija}
\index{razlika}
\index{podskup}
\index{univerzalni skup}
\index{komplement}

\key{Skup} je kolekcija elemenata.
Na primjer, skup
\[X=\{2,4,7\}\]
sadrži elemente 2, 4 i 7.
Simbol $\emptyset$ označava prazan skup,
a $|S|$ označava veličinu skupa $S$,
tj. broj elemenata u skupu.
Na primjer, u gornjem skupu je $|X|=3$.

Ako skup $S$ sadrži element $x$,
pišemo $x \in S$,
a inače pišemo $x \notin S$.
Na primjer, u gornjem skupu
\[4 \in X \hspace{10px}\textrm{i}\hspace{10px} 5 \notin X.\]

\begin{samepage}
Novi skupovi mogu se konstruirati pomoću operacija nad skupovima:
\begin{itemize}
\item \key{Presjek} $A \cap B$ sastoji se od elemenata
koji su i u $A$ i u $B$.
Na primjer, ako je $A=\{1,2,5\}$ i $B=\{2,4\}$,
tada je $A \cap B = \{2\}$.
\item \key{Unija} $A \cup B$ sastoji se od elemenata
koji su u $A$ ili u $B$ ili u oba.
Na primjer, ako je $A=\{3,7\}$ i $B=\{2,3,8\}$,
tada je $A \cup B = \{2,3,7,8\}$.
\item \key{Komplement} $\bar A$ sastoji se od elemenata
koji nisu u $A$.
Interpretacija komplementa ovisi o
\key{univerzalnom skupu}, koji sadrži sve moguće elemente.
Na primjer, ako je $A=\{1,2,5,7\}$ a univerzalni skup je
$\{1,2,\ldots,10\}$, tada je $\bar A = \{3,4,6,8,9,10\}$.
\item \key{Razlika} $A \setminus B = A \cap \bar B$
sastoji se od elemenata koji su u $A$, ali nisu u $B$.
Uočimo da $B$ može sadržavati elemente koji nisu u $A$.
Na primjer, ako je $A=\{2,3,7,8\}$ i $B=\{3,5,8\}$,
tada je $A \setminus B = \{2,7\}$.
\end{itemize}
\end{samepage}

Ako svaki element skupa $A$ pripada i skupu $S$,
kažemo da je $A$ \key{podskup} skupa $S$,
što označavamo s $A \subset S$.
Skup $S$ uvijek ima $2^{|S|}$ podskupova,
uključujući prazan skup.
Na primjer, podskupovi skupa $\{2,4,7\}$ su
\begin{center}
$\emptyset$,
$\{2\}$, $\{4\}$, $\{7\}$, $\{2,4\}$, $\{2,7\}$, $\{4,7\}$ i $\{2,4,7\}$.
\end{center}

Neki često korišteni skupovi su
$\mathbb{N}$ (prirodni brojevi),
$\mathbb{Z}$ (cijeli brojevi),
$\mathbb{Q}$ (racionalni brojevi) i
$\mathbb{R}$ (realni brojevi).
Skup $\mathbb{N}$
može se definirati na dva načina, ovisno
o situaciji:
bilo $\mathbb{N}=\{0,1,2,\ldots\}$
bilo $\mathbb{N}=\{1,2,3,...\}$.

Skup možemo konstruirati i pomoću pravila oblika
\[\{f(n) : n \in S\},\]
gdje je $f(n)$ neka funkcija.
Taj skup sadrži sve elemente oblika $f(n)$,
gdje je $n$ element skupa $S$.
Na primjer, skup
\[X=\{2n : n \in \mathbb{Z}\}\]
sadrži sve parne cijele brojeve.

\subsubsection{Logika}

\index{logika}
\index{negacija}
\index{konjunkcija}
\index{disjunkcija}
\index{implikacija}
\index{ekvivalencija}

Vrijednost logičkog izraza je
\key{istina} (1) ili \key{laž} (0).
Najvažniji logički operatori su
$\lnot$ (\key{negacija}),
$\land$ (\key{konjunkcija}),
$\lor$ (\key{disjunkcija}),
$\Rightarrow$ (\key{implikacija}) i
$\Leftrightarrow$ (\key{ekvivalencija}).
Sljedeća tablica prikazuje značenja tih operatora:

\begin{center}
\begin{tabular}{rr|rrrrrrr}
$A$ & $B$ & $\lnot A$ & $\lnot B$ & $A \land B$ & $A \lor B$ & $A \Rightarrow B$ & $A \Leftrightarrow B$ \\
\hline
0 & 0 & 1 & 1 & 0 & 0 & 1 & 1 \\
0 & 1 & 1 & 0 & 0 & 1 & 1 & 0 \\
1 & 0 & 0 & 1 & 0 & 1 & 0 & 0 \\
1 & 1 & 0 & 0 & 1 & 1 & 1 & 1 \\
\end{tabular}
\end{center}

Izraz $\lnot A$ ima suprotnu vrijednost od $A$.
Izraz $A \land B$ je istinit ako su i $A$ i $B$
istiniti,
a izraz $A \lor B$ je istinit ako je $A$ ili $B$ ili su oba
istiniti.
Izraz $A \Rightarrow B$ je istinit
ako je, kad god je $A$ istinit, i $B$ istinit.
Izraz $A \Leftrightarrow B$ je istinit
ako su $A$ i $B$ oba istinita ili oba lažna.

\index{predikat}

\key{Predikat} je izraz koji je istinit ili lažan
ovisno o svojim parametrima.
Predikate obično označavamo velikim slovima.
Na primjer, možemo definirati predikat $P(x)$
koji je istinit točno kada je $x$ prost broj.
Prema toj definiciji, $P(7)$ je istinit, a $P(8)$ je lažan.

\index{kvantifikator}

\key{Kvantifikator} povezuje logički izraz
s elementima skupa.
Najvažniji kvantifikatori su
$\forall$ (\key{za svaki}) i $\exists$ (\key{postoji}).
Na primjer,
\[\forall x (\exists y (y < x))\]
znači da za svaki element $x$ u skupu
postoji element $y$ u skupu
takav da je $y$ manji od $x$.
To je istinito u skupu cijelih brojeva,
ali lažno u skupu prirodnih brojeva.

Koristeći gore opisanu notaciju,
možemo izraziti mnogo vrsta logičkih tvrdnji.
Na primjer,
\[\forall x ((x>1 \land \lnot P(x)) \Rightarrow (\exists a (\exists b (a > 1 \land b > 1 \land x = ab))))\]
znači da, ako je broj $x$ veći od 1
i nije prost broj,
tada postoje brojevi $a$ i $b$
koji su veći od $1$ i čiji je proizvod $x$.
Ta tvrdnja je istinita u skupu cijelih brojeva.

\subsubsection{Funkcije}

Funkcija $\lfloor x \rfloor$ zaokružuje broj $x$
prema dolje na cijeli broj, a funkcija
$\lceil x \rceil$ zaokružuje broj $x$
prema gore na cijeli broj. Na primjer,
\[ \lfloor 3/2 \rfloor = 1 \hspace{10px} \textrm{i} \hspace{10px} \lceil 3/2 \rceil = 2.\]

Funkcije $\min(x_1,x_2,\ldots,x_n)$
i $\max(x_1,x_2,\ldots,x_n)$
daju najmanju i najveću od vrijednosti
$x_1,x_2,\ldots,x_n$.
Na primjer,
\[ \min(1,2,3)=1 \hspace{10px} \textrm{i} \hspace{10px} \max(1,2,3)=3.\]

\index{faktorijela}

\key{Faktorijela} $n!$ može se definirati kao
\[\prod_{x=1}^n x = 1 \cdot 2 \cdot 3 \cdot \ldots \cdot n\]
ili rekurzivno
\[
\begin{array}{lcl}
0! & = & 1 \\
n! & = & n \cdot (n-1)! \\
\end{array}
\]

\index{Fibonaccijev broj}

\key{Fibonaccijevi brojevi}
%\footnote{Fibonacci (c. 1175--1250) was an Italian mathematician.}
pojavljuju se u mnogim situacijama.
Mogu se definirati rekurzivno na sljedeći način:
\[
\begin{array}{lcl}
f(0) & = & 0 \\
f(1) & = & 1 \\
f(n) & = & f(n-1)+f(n-2) \\
\end{array}
\]
Prvi Fibonaccijevi brojevi su
\[0, 1, 1, 2, 3, 5, 8, 13, 21, 34, 55, \ldots\]
Postoji i zatvorena formula
za računanje Fibonaccijevih brojeva, koja se ponekad naziva
\index{Binetova formula} \key{Binetova formula}:
\[f(n)=\frac{(1 + \sqrt{5})^n - (1-\sqrt{5})^n}{2^n \sqrt{5}}.\]

\subsubsection{Logaritmi}

\index{logaritam}

\key{Logaritam} broja $x$
označava se s $\log_k(x)$, gdje je $k$ baza
logaritma.
Prema definiciji,
$\log_k(x)=a$ točno kada je $k^a=x$.

Korisno svojstvo logaritama jest
da je $\log_k(x)$ jednak broju puta
koliko moramo podijeliti $x$ s $k$ prije nego što dođemo 
do broja 1.
Na primjer, $\log_2(32)=5$
jer je potrebno 5 dijeljenja s 2:

\[32 \rightarrow 16 \rightarrow 8 \rightarrow 4 \rightarrow 2 \rightarrow 1 \]

Logaritmi se često koriste u analizi
algoritama, jer mnogi efikasni algoritmi
u svakom koraku nešto prepolove.
Zato efikasnost takvih algoritama možemo
procijeniti pomoću logaritama.

Logaritam proizvoda je
\[\log_k(ab) = \log_k(a)+\log_k(b),\]
a posljedično,
\[\log_k(x^n) = n \cdot \log_k(x).\]
Osim toga, logaritam kvocijenta je
\[\log_k\Big(\frac{a}{b}\Big) = \log_k(a)-\log_k(b).\]
Još jedna korisna formula je
\[\log_u(x) = \frac{\log_k(x)}{\log_k(u)},\]
a pomoću nje moguće je izračunati
logaritme po bilo kojoj bazi ako postoji način da
izračunamo logaritme po nekoj fiksnoj bazi.

\index{prirodni logaritam}

\key{Prirodni logaritam} $\ln(x)$ broja $x$
logaritam je kojemu je baza $e \approx 2.71828$.
Još jedno svojstvo logaritama jest da je
broj znamenki cijelog broja $x$ u bazi $b$ jednak
$\lfloor \log_b(x)+1 \rfloor$.
Na primjer, prikaz broja
$123$ u bazi $2$ je 1111011 i
$\lfloor \log_2(123)+1 \rfloor = 7$.

\section{Natjecanja i izvori}

\subsubsection{IOI}

Međunarodna informatička olimpijada (IOI)
godišnje je programersko natjecanje za
učenike srednjih škola.
Svaka zemlja smije poslati ekipu od
četiri učenika na natjecanje.
Obično sudjeluje oko 300 natjecatelja
iz 80 zemalja.

IOI se sastoji od dva petosatna natjecanja.
Na oba natjecanja od sudionika se traži da
riješe tri algoritamska zadatka različite težine.
Zadaci su podijeljeni na podzadatke,
od kojih svaki nosi određeni broj bodova.
Iako su natjecatelji podijeljeni u ekipe,
natječu se pojedinačno.

IOI-jev nastavni plan \cite{iois} propisuje teme
koje se mogu pojaviti u IOI zadacima.
Gotovo sve teme iz IOI-jeva nastavnog plana
obrađene su u ovoj knjizi.

Sudionici IOI-ja odabiru se putem
državnih natjecanja.
Prije IOI-ja organiziraju se mnoga regionalna natjecanja,
kao što su Baltička informatička olimpijada (BOI),
Srednjoeuropska informatička olimpijada (CEOI)
i Azijsko-pacifička informatička olimpijada (APIO).

Neke zemlje organiziraju online natjecanja za vježbu
za buduće sudionike IOI-ja,
kao što su Hrvatsko otvoreno natjecanje u informatici \cite{coci}
i Američka računalna olimpijada \cite{usaco}.
Osim toga, velika zbirka zadataka s poljskih natjecanja
dostupna je online \cite{main}.

\subsubsection{ICPC}

Međunarodno studentsko programersko natjecanje (ICPC)
godišnje je programersko natjecanje za studente.
Svaka ekipa na natjecanju sastoji se od tri studenta,
a za razliku od IOI-ja, studenti rade zajedno;
svakoj ekipi dostupno je samo jedno računalo.

ICPC se sastoji od nekoliko faza, a na kraju su
najbolje ekipe pozvane na Svjetsko finale.
Iako na natjecanju sudjeluju deseci tisuća
natjecatelja, dostupan je samo mali broj\footnote{Točan broj mjesta
u finalu varira od godine do godine; 2017. bilo je 133 mjesta u finalu.} mjesta u finalu,
pa je već i prolaz u finale
velik uspjeh u nekim regijama.

Na svakom ICPC natjecanju ekipe imaju pet sati vremena da
riješe oko deset algoritamskih zadataka.
Rješenje zadatka prihvaća se samo ako efikasno rješava
sve testne primjere.
Tijekom natjecanja natjecatelji mogu vidjeti rezultate drugih ekipa,
ali posljednji sat ljestvica je zamrznuta i nije
moguće vidjeti rezultate posljednjih predanih rješenja.

Teme koje se mogu pojaviti na ICPC-u nisu tako dobro
određene kao one na IOI-ju.
U svakom slučaju, jasno je da je na ICPC-u potrebno
više znanja, posebno više matematičkih vještina.

\subsubsection{Online natjecanja}

Postoje i mnoga online natjecanja otvorena za svakoga.
Trenutno je najaktivnija natjecateljska stranica Codeforces,
koja organizira natjecanja približno svaki tjedan.
Na Codeforcesu sudionici su podijeljeni u dvije divizije:
početnici se natječu u Div2, a iskusniji programeri u Div1.
Ostale natjecateljske stranice uključuju AtCoder, CS Academy, HackerRank i Topcoder.

Neke kompanije organiziraju online natjecanja s finalima na licu mjesta.
Primjeri takvih natjecanja su Facebook Hacker Cup,
Google Code Jam i Yandex.Algorithm.
Naravno, kompanije ta natjecanja koriste i za zapošljavanje:
dobar rezultat na natjecanju dobar je način da netko dokaže svoje vještine.

\subsubsection{Knjige}

Već postoje neke knjige (osim ove knjige) koje se
usredotočuju na natjecateljsko programiranje i algoritamsko rješavanje zadataka:

\begin{itemize}
\item S. S. Skiena i M. A. Revilla:
\emph{Programming Challenges: The Programming Contest Training Manual} \cite{ski03}
\item S. Halim i F. Halim:
\emph{Competitive Programming 3: The New Lower Bound of Programming Contests} \cite{hal13}
\item K. Diks i dr.: \emph{Looking for a Challenge? The Ultimate Problem Set from
the University of Warsaw Programming Competitions} \cite{dik12}
\end{itemize}

Prve dvije knjige namijenjene su početnicima,
dok posljednja knjiga sadrži napredniji materijal.

Naravno, i opće knjige o algoritmima prikladne su za
natjecatelje u programiranju.
Neke popularne knjige su:

\begin{itemize}
\item T. H. Cormen, C. E. Leiserson, R. L. Rivest i C. Stein:
\emph{Introduction to Algorithms} \cite{cor09}
\item J. Kleinberg i É. Tardos:
\emph{Algorithm Design} \cite{kle05}
\item S. S. Skiena:
\emph{The Algorithm Design Manual} \cite{ski08}
\end{itemize}
